\documentclass{article}
\usepackage{amsmath}
\usepackage{geometry}% 1-inch margins
\geometry{letterpaper, margin=1in}% 1-inch margins
\usepackage{titlesec}  % Allows customization of section/subsection titles
\titleformat{\section}[block]{\normalfont\Large\bfseries}{\thesection}{1em}{}
\titleformat{\subsection}[block]{\normalfont\large\bfseries}{\thesubsection}{1em}{}

\begin{document}
% Horizontal line
\hrule
\vspace{0.2cm}

% Centered title
\begin{center}
    \LARGE \textbf{EMCH 290} \\
    \Large Mechanical Engineering Thermodynamics \\
    \large Exam \#1 \textbf{SOLUTIONS}
\end{center}

% Horizontal line
\hrule
\vspace{0.2cm}

% Author and points possible
\begin{center}
    \textbf{Author:} Jacob Vaught \\
    \vspace{0.2cm}
    \textbf{Semester:} Fall 2023 \\
    \vspace{0.2cm}
    \textbf{Points Possible: 85 [21\% of FINAL Class Grade]}
\end{center}

% Another horizontal line
\hrule
\vspace{0.5cm}
\newpage
\section*{Problem 1[15 Points]}
Ground water has leaked into the tank to the depth as shown in the figure. The densities of the water and oil are, respectively, \(62 \, \text{lbm/ft}^3\) and \(55 \, \text{lbm/ft}^3\). The tank has 5 ft of water at the bottom and 30 ft of oil on top of the water.

\begin{enumerate}
    \item (5 Points)Calculate the pressure \(P_{\text{oil}}\) at the oil/water interface (\(h_{\text{oil}} = 30 \, \text{ft}\)) in \( \text{lbf/in}^2 \) (gage).
    \item (5 Points)Calculate the pressure \(P_{\text{water}}\) by water (\(h_{\text{water}} = 5 \, \text{ft}\)) in \( \text{lbf/in}^2 \) (gage).
    \item (2 Points)Calculate the total pressure on the bottom of the tank in \( \text{lbf/in}^2 \) (gage).
    \item (3 Points)What is the absolute pressure on the bottom of the tank at the above condition in \( \text{lbf/in}^2 \)? The atmospheric pressure is \(14.7 \, \text{lbf/in}^2\).
\end{enumerate}

\subsection*{Solutions}

\subsubsection*{Pressure at the Oil/Water Interface (P oil)}
To calculate the pressure \(P_{\text{oil}}\) at the oil/water interface, we use the formula:

\[
P = \rho \cdot g \cdot h
\]

Where \( \rho = 55 \, \text{lbm/ft}^3 \), and \( h = 30 \, \text{ft} \).

\[
P_{\text{oil}} = 55 \, \text{lbm/ft}^3 \times 30 \, \text{ft} = 1650 \, \text{lbf/ft}^2
\]

To convert this to \( \text{lbf/in}^2 \), we use the conversion \(1 \, \text{ft}^2 = 144 \, \text{in}^2\).

\[
P_{\text{oil}} = \frac{1650}{144} = 11.458 \, \text{lbf/in}^2
\]

\subsubsection*{Pressure by Water (P water)}
Similarly, for water:

\[
P_{\text{water}} = 62 \, \text{lbm/ft}^3 \times 5 \, \text{ft} = 9972 \, \text{lbf/ft}^2
\]

\[
P_{\text{water}} = \frac{310}{144} = 2.152 \, \text{lbf/in}^2
\]

\subsubsection*{Total Pressure on the Bottom of the Tank}
The total pressure \( P_{\text{total}} \) on the bottom of the tank is:

\[
P_{\text{total}} = P_{\text{oil}} + P_{\text{water}} = 2.152 \, \text{lbf/in}^2 + 11.458 \, \text{lbf/in}^2 = 13.61 \, \text{lbf/in}^2
\]

\subsubsection*{Absolute Pressure on the Bottom of the Tank}
The absolute pressure \( P_{\text{absolute}} \) is:

\[
P_{\text{absolute}} = P_{\text{total}} + P_{\text{atm}} = 13.61 \, \text{lbf/in}^2 + 14.7 \, \text{lbf/in}^2 = 28.31 \, \text{lbf/in}^2
\]
\newpage
\section*{Problem 2[20 Points]}
A gas is contained within a piston-cylinder assembly. The gas undergoes a process for which the pressure-volume relationship and the work equation are given as:
\[
PV^{1.5} = \text{Constant}
\]
\[
W = \frac{(P_2V_2 - P_1V_1)}{(n-1)}
\]
The initial pressure is \(2 \, \text{bar}\), the initial volume is \(0.1 \, \text{liter}\), and the final volume is \(0.2 \, \text{liter}\). The mass of the gas is \(5 \, \text{kg}\). The change in specific internal energy of the gas in the process \(\Delta u\) is \(40 \, \text{kJ/kg}\). If there are no significant changes in kinetic energy or potential energy, determine the heat transfer for the process.

\subsection*{Solutions}

\subsubsection*{Find the Final Pressure \(P_2\)}
Using the pressure-volume relationship \(PV^{1.5} = \text{Constant}\), we have:
\[
P_1V_1^{1.5} = P_2V_2^{1.5}
\]
Substitute \(P_1 = 2 \, \text{bar}\), \(V_1 = 0.1 \, \text{liter}\), and \(V_2 = 0.2 \, \text{liter}\):
\[
2 \times 0.1^{1.5} = P_2 \times 0.2^{1.5}
\]
Solving for \(P_2\):
\[
P_2 = \frac{2 \times 0.1^{1.5}}{0.2^{1.5}} = 0.7071 \, \text{bar}
\]

\subsubsection*{Find the Work Done \(W\)}
Using the work equation
\[
W = \frac{(P_2V_2 - P_1V_1)}{(n-1)}
\]
Since \(n = 1.5\), \(W\) becomes:
\[
W = \frac{(0.7071 \times 0.2 - 2 \times 0.1)}{(1.5-1)} = \frac{0.14142 - 0.2}{0.5} = -11.716 J
\]

\subsubsection*{Find the Heat Transfer \(Q\)}
The first law of thermodynamics for a closed system is:
\[
\Delta U = Q - W
\]
Given \(\Delta u = 40 \, \text{kJ/kg}\) and \(m = 5 \, \text{kg}\), the total internal energy change \(\Delta U\) is:
\[
\Delta U = 5 \times 40 = 200 \, \text{kJ}
\]
Since \(W = -0.011716 kJ\), \(Q\) is equal to \(\Delta U\):
\[
Q = 199.988 \, \text{kJ}
\]
\newpage
\section*{Problem 3[20 Points]}
A known gas within a piston-cylinder assembly undergoes a constant pressure expansion process. The piston has a face area of \(0.10 \, \text{m}^2\) and a mass of \(50 \, \text{kg}\). The initial volume is \(0.22 \, \text{m}^3\) and the final volume is \(0.34 \, \text{m}^3\). The piston-cylinder assembly is heated through the base. The change in specific internal energy of the gas is \(0.8 \, \text{kJ/kg}\) and the mass of gas is \(0.50 \, \text{kg}\). The piston and the cylinder walls are well insulated, and the piston moves smoothly in the cylinder. The atmospheric pressure is \(100 \, \text{kPa}\). Assuming that the friction between piston and cylinder is negligible and the acceleration of gravity is \(9.8 \, \text{m/s}^2\).

\begin{enumerate}
    \item (6 Points)The pressure by the gas in the cylinder in kPa.
    \item (6 Points)Calculate the work by \(P_{\text{gas}}\) in kJ.
    \item (3 Points)What is the change in internal energy of gas in kJ?
    \item (5 Points)Calculate the heat transfer \(Q\) from the hot plate to the gas in the cylinder in kJ.
\end{enumerate}

\subsection*{Solutions}

\subsubsection*{Pressure by the Gas in Cylinder \(P_{\text{gas}}\)}
To find the pressure exerted by the gas, we consider the force due to the weight of the piston and the atmospheric pressure:
\[
P_{\text{gas}} = P_{\text{atm}} + \frac{\text{Weight of Piston}}{\text{Area}}
\]
\[
P_{\text{gas}} = 100 \, \text{kPa} + \frac{(50 \, \text{kg} \times 9.8 \, \text{m/s}^2)}{0.10 \, \text{m}^2}
\]
\[
P_{\text{gas}} = 100 \, \text{kPa} + 4900 \, \text{Pa} = 100 \, \text{kPa} + 4.9 \, \text{kPa} = 104.9 \, \text{kPa}
\]

\subsubsection*{Work by \(P_{\text{gas}}\)}
The work \(W\) done during a constant pressure process is:
\[
W = P \times \Delta V
\]
\[
W = 104.9 \, \text{kPa} \times (0.34 \, \text{m}^3 - 0.22 \, \text{m}^3) = 104.9 \, \text{kPa} \times 0.12 \, \text{m}^3 = 12.588 \, \text{kJ}
\]

\subsubsection*{Change in Internal Energy of Gas}
The change in internal energy \( \Delta U \) is given as \(0.8 \, \text{kJ/kg}\) for \(0.50 \, \text{kg}\) of gas:
\[
\Delta U = 0.8 \, \text{kJ/kg} \times 0.50 \, \text{kg} = 0.4 \, \text{kJ}
\]

\subsubsection*{Heat Transfer \(Q\)}
By the first law of thermodynamics,
\[
Q = \Delta U + W
\]
\[
Q = 0.4 \, \text{kJ} + 12.588 \, \text{kJ} = 12.988 \, \text{kJ}
\]


\newpage
\section*{Problem 4[20 Points]}
A Condor weighing \(35 \, \text{lbf}\) fell from a height of \(15,000 \, \text{ft}\) above the surface of a lake with an initial velocity of \(50 \, \text{mph}\) during a diving event (free falling). For \(g = 30.15 \, \text{ft/s}^2\), determine:
\begin{enumerate}
    \item (6 Points) The mass of the Condor, in \(\text{lbm}\).
    \item (6 Points) The change in gravitational potential energy of the Condor \(\Delta PE\), in \(\text{ft} \cdot \text{lbf}\).
    \item (8 Points) If the final velocity of the Condor is \(20 \, \text{ft/s}\) before falling into water, what is the change in kinetic energy of the Condor \(\Delta KE\), in \(\text{ft-lbf}\)?
\end{enumerate}

\subsection*{Solutions}

\subsubsection*{Mass of the Condor in lbm}
Using the weight formula \(W = mg\), we find \(m\) as follows. Given that \( \text{lbf} = 32.174 \, \text{lbm} \times \text{ft/s}^2 \), we have:
\[
m = \frac{W}{g} = \frac{35 \, \text{lbf} \times 32.174 \, \text{lbm ft/s}^2/\text{lbf}}{30.15 \, \text{ft/s}^2} = 37.35 \, \text{lbm}
\]


\subsubsection*{Change in Gravitational Potential Energy \(\Delta PE\)}
The change in potential energy is given by \(\Delta PE = mgh\). Including the conversion factor, we have:
\[
\Delta PE = 37.35 \, \text{lbm} \times 30.15 \, \text{ft/s}^2 \times 15,000 \, \text{ft} \times \frac{1}{32.174 \, \text{lbf/(lbm} \, \text{ft/s}^2)} 
\]
\[
= 16,891,537.5 \, \text{ft} \cdot \text{lbf} \times \frac{1}{32.174 \, \text{lbf/(lbm} \, \text{ft/s}^2)} = 525,005.8 \, \text{ft} \cdot \text{lbf}
\]


\subsubsection*{Change in Kinetic Energy \(\Delta KE\)}
The initial kinetic energy \(KE_i\) is \(\frac{1}{2} mv^2\). The initial velocity is \(50 \, \text{mph} = 73.33 \, \text{ft/s}\):
\[
KE_i = \frac{1}{2} \times 37.35 \, \text{lbm} \times (73.33 \, \text{ft/s})^2 = 100,429.9 \, \text{ft} \cdot \text{lbm}
\]
The final kinetic energy \(KE_f\) with a final velocity of \(20 \, \text{ft/s}\) is:
\[
KE_f = \frac{1}{2} \times 37.35 \, \text{lbm} \times (20 \, \text{ft/s})^2 = 7,470 \, \text{ft} \cdot \text{lbm}
\]
The change in kinetic energy \(\Delta KE\) is \(KE_f - KE_i\):
\[
\Delta KE = 7,470 \, \text{ft} \cdot \text{lbm} - 100,429.9 \, \text{ft} \cdot \text{lbm} = -92959.9 \, \text{ft} \cdot \text{lbm}  \times \frac{1}{32.174 \, \text{lbf/(lbm} \, \text{ft/s}^2)} = 2889.28 \, \text{ft} \cdot \text{lbf}
\]



\newpage
\section*{Problem 5[4 Points]}
A refrigeration cycle operating as shown in the figure has a coefficient of performance \( \beta = 3.0 \). For the cycle, \( Q_{\text{out}} = 300 \, \text{kJ} \).

\begin{enumerate}
    \item (2 Points)Determine \( Q_{\text{in}} \) in kJ.
    \item (2 Points)Determine \( W_{\text{cycle}} \) in kJ.
\end{enumerate}

\subsection*{Solutions}
\subsubsection*{(a) Determine \( Q_{\text{in}} \) in kJ}
The coefficient of performance \( \beta \) is given by:
\[
\beta = \frac{Q_{\text{out}}-W_{\text{cycle}}}{W_{\text{cycle}}}
\]
We can rearrange this equation to find \( W_{\text{cycle}} \):
\[
W_{\text{cycle}} = \frac{Q_{\text{out}}}{\beta + 1} = \frac{300 \, \text{kJ}}{4.0} = 75 \, \text{kJ}
\]
Now, the energy balance for the cycle can be written as:
\[
Q_{\text{in}} = Q_{\text{out}} - W_{\text{cycle}}
\]
\[
Q_{\text{in}} = 300 \, \text{kJ} - 75 \, \text{kJ} = 225 \, \text{kJ}
\]

\subsubsection*{(b) Determine \( W_{\text{cycle}} \) in kJ}
From part (a), we found \( W_{\text{cycle}} = 75 \, \text{kJ} \).

\newpage
\section*{Problem 6[6 Points]}
Compare a control mass and control volume system. Present example for each to support your answer.
\subsection*{Solutions}

\subsubsection*{Control Mass vs. Control Volume Systems}
A control mass system refers to a fixed quantity of matter, and no mass crosses the boundary of the system. A control volume system allows mass to cross its boundaries.
\newline \newline
\textbf{Example for Control Mass}: A piston-cylinder arrangement with a closed end, where the air inside cannot escape, is an example of a control mass system.
\newline \newline
\textbf{Example for Control Volume}: A wind tunnel where air flows through it is an example of a control volume system.


\end{document}
